% !PRINTED 142
\weberpage{142}
%<<CONT

\weq{6}{\begin{aligned}
\frac{H(v)}{\Theta(v)} &= \sqrt{\varkappa}\operatorname{sn}v,\\
\frac{H(v+K)}{\Theta(v)} &= \sqrt{\frac{\varkappa}{\varkappa'}}\operatorname{cn}v,\\
\frac{\Theta(v+K)}{\Theta(v)} &= \frac{1}{\sqrt{\varkappa'}}\operatorname{dn}v.
\end{aligned}}
Die Funktionen $\Theta$, $H$ haben, wie aus (3) leicht folgt, die
durch die folgenden Gleichungen ausgedrückte Periodizität:
\weq{7}{\begin{aligned}
\Theta(v+2nK) &= \Theta(v),\\
\Theta(v+2niK') &= (-1)^n e^{\frac{n^2\pi K'}{K}-\frac{\pi inv}{K}}\Theta(v),
\end{aligned}}
\weq{8}{\begin{aligned}
H(v+2nK) &= (-1)^n H(v),\\
H(v+2niK') &= (-1)^n e^{\frac{n^2\pi K'}{K}-\frac{\pi inv}{K}}H(v).
\end{aligned}}
Es ist bisweilen nützlich, die Gleichungen (3), (7), (8)
folgendermaßen zusammenzufassen:
\weq{9}{\Phi(v+niK') = i^n e^{\frac{n^2\pi i}{4}\frac{K'}{K}-\frac{\pi inv}{2K}}\Psi(v),}
worin, wenn $n$ \gesperrt{gerade}, $\Phi$ und $\Psi$ beide gleich $\Theta$ oder beide
gleich $H$, wenn $n$ ungerade, $\Phi=\Theta$, $\Psi=H$ oder $\Phi=H$,
$\Psi=\Theta$ zu setzen ist.

Aus \S\,22 (2), (4) ergeben sich die Additionsformeln:
\weq{10}{\begin{aligned}
\Theta^2(0)\Theta(u+v)\Theta(u-v) &= \Theta^2(u)\Theta^2(v) - H^2(u)H^2(v),\\
\Theta^2(0)H(u+v)H(u-v) &= H^2(u)\Theta^2(v) - \Theta^2(u)H^2(v),
\end{aligned}}
wofür man auch schreiben kann:
\weq{11}{\begin{aligned}
\Theta^2(0)\Theta(u+v)\Theta(u-v) &= \Theta^2(u)\Theta^2(v)(1-\varkappa^2\operatorname{sn}^2u\operatorname{sn}^2v),\\
\Theta^2(0)H(u+v)H(u-v) &= \Theta^2(u)\Theta^2(v)\varkappa(\operatorname{sn}^2u-\operatorname{sn}^2v).
\end{aligned}}
Wir gehen nun dazu über, die Sätze über die $\vartheta$-Funktionen
auf die elliptischen Funktionen zu übertragen und beginnen mit
dem Additionstheorem.

\begin{center}
\textbf{\S\,44. \gesperrt{Additionstheorem der elliptischen Funktionen}.}
\end{center}

Aus den vier $\vartheta$-Funktionen lassen sich im ganzen zwölf
Quotienten bilden, die nach \S\,42 durch elliptische Funktionen
ausdrückbar sind. Bildet man diese Quotienten für das Argument
$u+v$, so kann man diese nach \S\,22 durch $\vartheta$-Funktionen von $u$
und von $v$ einzeln, also auch durch die entsprechenden elliptischen
 Funktionen darstellen, und zwar immer so, daß der Nenner
%CONT>>
